\documentclass{amsart}
\usepackage{amsmath,amssymb,amsthm,hyperref}
\usepackage{bm}

\theoremstyle{definition}
\newtheorem{definition}{Definition}
\newtheorem{assumption}{Assumption}
\newtheorem{model}{Model}
\theoremstyle{plain}
\newtheorem{proposition}{Proposition}
\newtheorem{theorem}{Theorem}
\newtheorem{corollary}{Corollary}
\newtheorem{lemma}{Lemma}

\begin{document}

\title{An Integrated Structural--Interpersonal Framework for Team Learning Behavior}
\maketitle

\begin{abstract}
We formulate an integrated framework that jointly models team \emph{structural features}
and \emph{shared interpersonal beliefs} (psychological safety) as determinants of team
\emph{learning behavior} and downstream \emph{performance}. The framework is stated as a
recursive structural model on a directed acyclic graph (DAG). We make four claims precise
and prove them: (i) the structural features act partly as \emph{antecedents} of learning
behavior \emph{mediated} by psychological safety; (ii) psychological safety also acts as a
\emph{moderator} of the structure--learning relationship, so the two classes of variables
are \emph{not} independent additive predictors; (iii) under stated assumptions the
direct, indirect (mediated), and moderated effects are identified and the total effect
decomposes exactly into them; and (iv) the framework is simultaneously testable by
quantitative (parametric estimation of the structural coefficients) and qualitative
(mechanism-tracing of the proposed paths) methods. We prove an exact effect-decomposition
theorem, an identification theorem, and a non-reducibility theorem showing that any purely
additive ``independent predictors'' model is strictly misspecified whenever a genuine
interaction is present.
\end{abstract}

\section{Formalization}

We treat the substantive question --- how structural features and shared interpersonal
beliefs jointly produce team learning --- as a problem of specifying and analyzing a
generative structural model over team-level random variables, in the sense of Wright's path
analysis and Pearl's structural causal models. The substantive content (the directions of
the arrows and the role of psychological safety as both mediator and moderator) is the
modeling hypothesis; the mathematical content is the proof that, under explicit and
checkable assumptions, this hypothesis yields identified, decomposable, and testable
quantities, and that it is strictly richer than the additive alternative it is meant to
replace.

\begin{definition}[Team variables]\label{def:vars}
Fix a population of work teams indexed by a probability space $(\Omega,\mathcal F,\mathbb P)$.
For a team $\omega$ define the following team-level random vectors/scalars, all assumed to
have finite second moments:
\begin{itemize}
\item $\bm X=(X_1,\dots,X_p)\in\mathbb R^p$: \emph{structural features}
(team design, leader coaching, context support, composition);
\item $S\in\mathbb R$: \emph{psychological safety}, the shared belief that the team is safe
for interpersonal risk-taking, operationalized as a team-level latent obtained by
aggregating member perceptions (see Definition~\ref{def:shared});
\item $L\in\mathbb R$: \emph{learning behavior} (feedback seeking, information sharing,
help seeking, experimentation, error discussion), a team-level composite;
\item $P\in\mathbb R$: downstream \emph{team performance}.
\end{itemize}
\end{definition}

\begin{definition}[Sharedness / aggregation validity]\label{def:shared}
Let $S_{ij}$ be member $j$'s perception in team $i$, $j=1,\dots,n_i$. The construct $S$ is
\emph{shared} if the within-team agreement is high and between-team variance dominates,
formalized by an intraclass correlation
\[
\mathrm{ICC}(1)=\frac{\sigma^2_{\text{between}}}{\sigma^2_{\text{between}}+\sigma^2_{\text{within}}}
\]
bounded away from $0$, in which case the team mean $S_i=\frac1{n_i}\sum_j S_{ij}$ is a
consistent estimator of the team-level latent as $n_i\to\infty$ and the between-team model
below is well posed. (This is the measurement-model precondition that licenses treating $S$
as a single team-level variable rather than $n_i$ individual variables.)
\end{definition}

\section{The integrated structural model}

\begin{model}[Integrated structural--interpersonal model]\label{mod:main}
The framework asserts the recursive (acyclic) generative process
\begin{align}
S &= \alpha_0 + \bm\alpha^\top \bm X + \varepsilon_S, \label{eq:S}\\
L &= \beta_0 + \bm\beta^\top \bm X + \gamma\, S
      + (\bm\delta^\top \bm X)\, S + \varepsilon_L, \label{eq:L}\\
P &= \eta_0 + \theta\, L + \bm\phi^\top \bm X + \varepsilon_P, \label{eq:P}
\end{align}
with the ordering $\bm X \prec S \prec L \prec P$ on the DAG $\bm X\to S\to L\to P$,
$\bm X\to L$, $\bm X\to P$. Here $\bm\alpha,\bm\beta,\bm\delta,\bm\phi\in\mathbb R^p$ and
$\alpha_0,\beta_0,\eta_0,\gamma,\theta\in\mathbb R$ are parameters, and
$(\varepsilon_S,\varepsilon_L,\varepsilon_P)$ are mean-zero disturbances.
The term $(\bm\delta^\top\bm X)\,S$ encodes \emph{moderation}: the slope of $L$ on $S$ is
$\gamma+\bm\delta^\top\bm X$, i.e.\ structure conditions the safety--learning link.
\end{model}

The two substantive roles are read directly off the equations:

\begin{itemize}
\item \textbf{Antecedent + mediator.} In \eqref{eq:S} structural features are
\emph{antecedents} of psychological safety; in \eqref{eq:L} safety $S$ transmits part of
the effect of $\bm X$ onto $L$. Thus $S$ \emph{mediates} the structure--learning link.
\item \textbf{Moderator.} The cross term $(\bm\delta^\top\bm X)S$ makes safety also a
\emph{moderator}: the marginal return to safety in producing learning depends on the
structural configuration (and symmetrically the return to structure depends on safety).
This is precisely the ``not independent predictors'' requirement: $\partial^2 L/\partial
\bm X\,\partial S=\bm\delta\neq\bm 0$ in general.
\end{itemize}

We use the following standard exogeneity assumption, which is the minimal substantive
commitment needed for the parameters to carry causal meaning; in observational team data it
is the hypothesis that design/coaching/support/composition are exogenous (or have been made
so by design, e.g.\ randomized coaching interventions) relative to the residual unobserved
team heterogeneity.

\begin{assumption}[Recursive exogeneity]\label{ass:exog}
Conditional on the regressors appearing in each equation, the disturbances are mean
independent and mutually orthogonal across equations:
\[
\mathbb E[\varepsilon_S\mid \bm X]=0,\quad
\mathbb E[\varepsilon_L\mid \bm X,S]=0,\quad
\mathbb E[\varepsilon_P\mid \bm X,L]=0,
\]
and $\mathbb E[\varepsilon_S\varepsilon_L]=\mathbb E[\varepsilon_S\varepsilon_P]
=\mathbb E[\varepsilon_L\varepsilon_P]=0$. We also assume
$\mathbb E[\varepsilon_L\mid \bm X,S]=0$ holds jointly with the interaction regressor,
i.e.\ $\mathbb E[\varepsilon_L\mid \bm X,S,(\bm\delta^\top\bm X)S]=0$, which follows from
the first display because $(\bm\delta^\top\bm X)S$ is a measurable function of $(\bm X,S)$.
\end{assumption}

\section{Effect decomposition (mediation + moderation)}

We now make precise the sense in which the total influence of structure on learning splits
into a direct path and a safety-mediated path, and how moderation enters. Write
$\bm x,\bm x'$ for two structural configurations and define the \emph{conditional total
effect} of moving $\bm X$ from $\bm x'$ to $\bm x$ on $L$ as the difference of conditional
means $\mathbb E[L\mid \bm X=\bm x]-\mathbb E[L\mid \bm X=\bm x']$ under
Model~\ref{mod:main}.

\begin{theorem}[Exact total-effect decomposition]\label{thm:decomp}
Under Model~\ref{mod:main} and Assumption~\ref{ass:exog}, for any $\bm x,\bm x'\in\mathbb
R^p$,
\begin{equation}\label{eq:decomp}
\underbrace{\mathbb E[L\mid \bm X=\bm x]-\mathbb E[L\mid \bm X=\bm x']}_{\text{total effect}}
=\underbrace{\bm\beta^\top(\bm x-\bm x')}_{\text{direct}}
+\underbrace{(\gamma+\bm\delta^\top\bm x)\,\bm\alpha^\top\bm x
   -(\gamma+\bm\delta^\top\bm x')\,\bm\alpha^\top\bm x'}_{\text{mediated through }S}
+\underbrace{(\bm\delta^\top\bm x-\bm\delta^\top\bm x')\,\alpha_0}_{\text{moderated baseline}}.
\end{equation}
In the special case $\bm\delta=\bm 0$ (no moderation) this collapses to the classical linear
mediation identity
\begin{equation}\label{eq:classmed}
\mathbb E[L\mid \bm X=\bm x]-\mathbb E[L\mid \bm X=\bm x']
=\bigl(\bm\beta+\gamma\,\bm\alpha\bigr)^\top(\bm x-\bm x'),
\end{equation}
i.e.\ \emph{total} $=$ \emph{direct} ($\bm\beta$) $+$ \emph{indirect} ($\gamma\bm\alpha$).
\end{theorem}

\begin{proof}
By Assumption~\ref{ass:exog}, $\mathbb E[\varepsilon_L\mid \bm X,S]=0$, so taking the
conditional expectation of \eqref{eq:L} given $\bm X=\bm x$ and then iterating expectations
over $S$,
\begin{align}
\mathbb E[L\mid \bm X=\bm x]
&=\beta_0+\bm\beta^\top\bm x+\gamma\,\mathbb E[S\mid \bm X=\bm x]
   +\bm\delta^\top\bm x\,\mathbb E[S\mid \bm X=\bm x],\label{eq:step1}
\end{align}
where we used that $(\bm\delta^\top\bm x)$ is constant given $\bm X=\bm x$, so
$\mathbb E[(\bm\delta^\top\bm X)S\mid \bm X=\bm x]=(\bm\delta^\top\bm x)\,
\mathbb E[S\mid \bm X=\bm x]$.
By \eqref{eq:S} and $\mathbb E[\varepsilon_S\mid \bm X]=0$,
\begin{equation}\label{eq:Smean}
\mathbb E[S\mid \bm X=\bm x]=\alpha_0+\bm\alpha^\top\bm x.
\end{equation}
Substituting \eqref{eq:Smean} into \eqref{eq:step1},
\begin{equation}\label{eq:Lmean}
\mathbb E[L\mid \bm X=\bm x]
=\beta_0+\bm\beta^\top\bm x+(\gamma+\bm\delta^\top\bm x)(\alpha_0+\bm\alpha^\top\bm x).
\end{equation}
Writing the same identity at $\bm x'$ and subtracting,
\begin{align*}
\mathbb E[L\mid \bm X=\bm x]-\mathbb E[L\mid \bm X=\bm x']
&=\bm\beta^\top(\bm x-\bm x')\\
&\quad+(\gamma+\bm\delta^\top\bm x)(\alpha_0+\bm\alpha^\top\bm x)
       -(\gamma+\bm\delta^\top\bm x')(\alpha_0+\bm\alpha^\top\bm x').
\end{align*}
Expanding the last two products and grouping the terms multiplying $\alpha_0$ separately
from those multiplying $\bm\alpha^\top(\cdot)$ gives exactly \eqref{eq:decomp}.
If $\bm\delta=\bm 0$, then \eqref{eq:Lmean} becomes
$\beta_0+\bm\beta^\top\bm x+\gamma(\alpha_0+\bm\alpha^\top\bm x)$, whose difference at
$\bm x,\bm x'$ is $(\bm\beta+\gamma\bm\alpha)^\top(\bm x-\bm x')$, which is
\eqref{eq:classmed}.
\end{proof}

\begin{corollary}[Performance pathway]\label{cor:perf}
Under Model~\ref{mod:main} and Assumption~\ref{ass:exog}, learning behavior is the channel
through which the structure--safety system reaches performance:
\[
\mathbb E[P\mid \bm X=\bm x]
=\eta_0+\bm\phi^\top\bm x+\theta\,\mathbb E[L\mid \bm X=\bm x],
\]
with $\mathbb E[L\mid \bm X=\bm x]$ given by \eqref{eq:Lmean}. Hence the total effect of
structure on performance equals the \emph{direct} structural part $\bm\phi^\top(\bm x-\bm
x')$ plus $\theta$ times the total effect on learning from Theorem~\ref{thm:decomp}.
\end{corollary}
\begin{proof}
Take the conditional expectation of \eqref{eq:P} given $\bm X=\bm x$; by
Assumption~\ref{ass:exog} $\mathbb E[\varepsilon_P\mid \bm X,L]=0$ and the law of iterated
expectations give $\mathbb E[P\mid\bm X=\bm x]=\eta_0+\bm\phi^\top\bm x
+\theta\,\mathbb E[L\mid\bm X=\bm x]$. Subtracting the analogous expression at $\bm x'$ and
applying Theorem~\ref{thm:decomp} to $\mathbb E[L\mid\bm X=\bm x]-\mathbb E[L\mid\bm
X=\bm x']$ proves the claim.
\end{proof}

\section{Identification}

The decomposition is only useful if its ingredients are recoverable from data. We give
sufficient conditions.

\begin{theorem}[Parameter identification]\label{thm:ident}
Suppose, in addition to Assumption~\ref{ass:exog}:
\begin{enumerate}
\item[(R1)] the second-moment matrix of the regressor vector
$\bm Z:=(1,\bm X^\top,S,(\bm X S)^\top)^\top\in\mathbb R^{2p+2}$ that appears in
\eqref{eq:L} is nonsingular, $\Sigma_Z:=\mathbb E[\bm Z\bm Z^\top]\succ 0$;
\item[(R2)] the second-moment matrix of $(1,\bm X^\top)$ is nonsingular;
\item[(R3)] the second-moment matrix of $(1,\bm X^\top,L)$ is nonsingular.
\end{enumerate}
Then all structural parameters
$(\alpha_0,\bm\alpha)$, $(\beta_0,\bm\beta,\gamma,\bm\delta)$, and
$(\eta_0,\bm\phi,\theta)$ are uniquely determined (identified) by the joint distribution of
$(\bm X,S,L,P)$, and consequently every term in \eqref{eq:decomp} and Corollary~\ref{cor:perf}
is identified.
\end{theorem}

\begin{proof}
\emph{Equation \eqref{eq:S}.} By Assumption~\ref{ass:exog}, $\mathbb E[\varepsilon_S\mid\bm
X]=0$ implies $\mathbb E[(1,\bm X^\top)^\top\varepsilon_S]=\bm0$. Hence
$(\alpha_0,\bm\alpha)$ satisfy the normal equations
$\mathbb E[(1,\bm X^\top)^\top(1,\bm X^\top)](\alpha_0,\bm\alpha^\top)^\top
=\mathbb E[(1,\bm X^\top)^\top S]$. By (R2) the coefficient matrix is invertible, so
$(\alpha_0,\bm\alpha)$ is the unique solution; it equals the population least-squares
projection of $S$ on $(1,\bm X)$.

\emph{Equation \eqref{eq:L}.} Assumption~\ref{ass:exog} gives
$\mathbb E[\bm Z\,\varepsilon_L]=\bm0$ because each component of $\bm Z$ is a measurable
function of $(\bm X,S)$ and $\mathbb E[\varepsilon_L\mid \bm X,S]=0$. Therefore the
parameter vector $\bm b:=(\beta_0,\bm\beta^\top,\gamma,\bm\delta^\top)^\top$ solves
$\Sigma_Z\,\bm b=\mathbb E[\bm Z\,L]$, and by (R1) $\Sigma_Z\succ0$ is invertible, so
$\bm b=\Sigma_Z^{-1}\mathbb E[\bm Z\,L]$ is unique. In particular $\bm\delta$ is identified,
so the presence/absence of moderation is empirically decidable.

\emph{Equation \eqref{eq:P}.} Identically, $\mathbb E[\varepsilon_P\mid\bm X,L]=0$ gives
$\mathbb E[(1,\bm X^\top,L)^\top\varepsilon_P]=\bm0$; by (R3) the corresponding
second-moment matrix is invertible, so $(\eta_0,\bm\phi,\theta)$ is the unique solution of
its normal equations.

Since every structural parameter is a fixed (matrix-algebraic) functional of population
moments of $(\bm X,S,L,P)$, and \eqref{eq:decomp} and Corollary~\ref{cor:perf} are
continuous functions of these parameters and of $\bm x,\bm x'$, each effect term is
identified.
\end{proof}

\begin{remark}[Estimation and inference]
Theorem~\ref{thm:ident} is constructive: it exhibits each parameter as a population
least-squares projection. Replacing population moments by sample moments yields the ordinary
least-squares / SEM estimators, which are consistent and asymptotically normal under (R1)--(R3)
and finite fourth moments, so all effects in \eqref{eq:decomp} admit standard (e.g.\
delta-method or bootstrap) confidence intervals. The indirect effect $\gamma\bm\alpha$ in the
no-moderation case is the classical product-of-coefficients mediation estimand.
\end{remark}

\section{Why a joint model: non-reducibility to independent predictors}

The problem demands a framework that models the two classes of variables \emph{jointly}
rather than as independent predictors. We prove that this is not a stylistic preference: when
moderation is genuinely present ($\bm\delta\neq\bm0$), \emph{every} additive model in
$(\bm X,S)$ is strictly misspecified, so the joint formulation is necessary, not optional.

\begin{theorem}[Strict superiority of the joint model]\label{thm:nonred}
Let the data be generated by Model~\ref{mod:main} with $\bm\delta\neq\bm0$ and with the
interaction regressor not collinear with the additive ones (formally, the residual of
$(\bm X S)$ after projecting it on $(1,\bm X,S)$ is nonzero in $L^2$; this holds under (R1)).
Let
\[
g_{\mathrm{add}}(\bm X,S)=\mu_0+\bm\mu^\top\bm X+\nu\,S
\]
range over all affine (additive, no-interaction) predictors, and let
\[
g_{\mathrm{int}}(\bm X,S)=\beta_0+\bm\beta^\top\bm X+\gamma S+(\bm\delta^\top\bm X)S
\]
be the conditional mean of $L$ under \eqref{eq:L}. Then
\[
\min_{\mu_0,\bm\mu,\nu}\ \mathbb E\bigl[(L-g_{\mathrm{add}}(\bm X,S))^2\bigr]
\;>\;
\mathbb E\bigl[(L-g_{\mathrm{int}}(\bm X,S))^2\bigr],
\]
with strict inequality. Equivalently, the best additive predictor of learning behavior from
structure and safety has strictly larger mean-squared error than the integrated model; the
gap equals the variance of the omitted interaction's $L^2$-projection.
\end{theorem}

\begin{proof}
Write $W:=(\bm\delta^\top\bm X)S$ and let $\hat W$ be its $L^2$ orthogonal projection onto
$\mathcal A:=\overline{\mathrm{span}}\{1,X_1,\dots,X_p,S\}$, with residual
$R:=W-\hat W$, so $R\perp\mathcal A$ and, by hypothesis, $\mathbb E[R^2]=\sigma_R^2>0$.

By \eqref{eq:L} and $\mathbb E[\varepsilon_L\mid\bm X,S]=0$ we have
$L=g_{\mathrm{int}}(\bm X,S)+\varepsilon_L$ with $\varepsilon_L\perp\mathcal A$ and also
$\varepsilon_L\perp W$ (since $W$ is $(\bm X,S)$-measurable), hence
$\varepsilon_L\perp R$. The best additive predictor is the $L^2$ projection of $L$ onto
$\mathcal A$. Decompose
\[
L=\bigl(\beta_0+\bm\beta^\top\bm X+\gamma S+\hat W\bigr)+R+\varepsilon_L,
\]
where the first bracket lies in $\mathcal A$, while $R\perp\mathcal A$ and
$\varepsilon_L\perp\mathcal A$. Therefore the projection of $L$ onto $\mathcal A$ is exactly
the first bracket, and the additive minimal MSE is
\[
\min_{\mathcal A}\mathbb E[(L-g_{\mathrm{add}})^2]
=\mathbb E[(R+\varepsilon_L)^2]
=\mathbb E[R^2]+\mathbb E[\varepsilon_L^2]
=\sigma_R^2+\mathbb E[\varepsilon_L^2],
\]
using $R\perp\varepsilon_L$. On the other hand,
\[
\mathbb E[(L-g_{\mathrm{int}}(\bm X,S))^2]=\mathbb E[\varepsilon_L^2].
\]
Subtracting, the additive optimum exceeds the integrated MSE by exactly
$\sigma_R^2>0$, which is strict because $\bm\delta\neq\bm0$ and the non-collinearity
hypothesis force $\sigma_R^2>0$.
\end{proof}

\begin{corollary}[Bias of the ``independent predictors'' coefficients]\label{cor:bias}
Under the hypotheses of Theorem~\ref{thm:nonred}, the additive least-squares slope on $S$,
\[
\nu^\ast=\arg\min_\nu \min_{\mu_0,\bm\mu}\mathbb E[(L-\mu_0-\bm\mu^\top\bm X-\nu S)^2],
\]
generally differs from $\gamma$ (the true safety slope at the structural mean), and from
$\gamma+\bm\delta^\top\bm x$ at any particular configuration $\bm x$. Hence treating safety
and structure as independent additive predictors not only loses fit
(Theorem~\ref{thm:nonred}) but yields a biased, configuration-blind estimate of the
safety--learning relationship.
\end{corollary}
\begin{proof}
$\nu^\ast$ is the coefficient on $S$ in the projection of $L$ onto $\mathcal A$, i.e.\ on
$g_{\mathrm{int}}$'s projection. Since the projection of $\hat W=$ proj$_{\mathcal A}
[(\bm\delta^\top\bm X)S]$ loads on $S$ whenever $\mathbb E[(\bm\delta^\top\bm X)S\,\cdot
S]\neq 0$ after partialling out $(1,\bm X)$, the $S$-coefficient absorbs part of the
interaction, so $\nu^\ast\neq\gamma$ in general. The true marginal effect of $S$ on $L$ at
$\bm X=\bm x$ is $\partial \mathbb E[L\mid\bm X=\bm x,S=s]/\partial s=\gamma+\bm\delta^\top
\bm x$ by \eqref{eq:L}, which varies with $\bm x$ and therefore cannot equal the single
constant $\nu^\ast$ unless $\bm\delta=\bm0$.
\end{proof}

\section{Dual quantitative--qualitative testability}

The framework is designed to be investigated by both methodological traditions; we state the
correspondence precisely.

\begin{proposition}[Quantitative test]\label{prop:quant}
Under Assumption~\ref{ass:exog} and (R1)--(R3), the following are point-identified and
estimable by least squares / SEM, each with a falsifiable null hypothesis:
\begin{enumerate}
\item \emph{Antecedent claim} $\bm\alpha\neq\bm0$ (structure predicts safety): test
$H_0:\bm\alpha=\bm0$ in \eqref{eq:S}.
\item \emph{Mediation claim} $\gamma\bm\alpha\neq\bm0$ (indirect path): test
$H_0:\gamma\bm\alpha=\bm0$ via the product-of-coefficients / bootstrap test.
\item \emph{Moderation claim} $\bm\delta\neq\bm0$: test the interaction block in
\eqref{eq:L}; by Theorem~\ref{thm:nonred} a significant $\bm\delta$ is equivalent to a
significant improvement in fit of the integrated over the additive model.
\item \emph{Performance channel} $\theta\neq0$ and the indirect effect of structure on $P$
through $L$ (Corollary~\ref{cor:perf}): test $H_0:\theta=0$ and the corresponding mediated
effect.
\end{enumerate}
\end{proposition}
\begin{proof}
Each parameter is identified by Theorem~\ref{thm:ident} and consistently estimable by the
associated least-squares projection (Remark after Theorem~\ref{thm:ident}); a parameter
restriction is a linear (or smooth) hypothesis on those estimates, hence falsifiable by a
Wald / likelihood-ratio / bootstrap test. The equivalence in item~3 is
Theorem~\ref{thm:nonred}: $\bm\delta\neq\bm0$ iff the additive model has strictly larger
population MSE.
\end{proof}

\begin{proposition}[Qualitative / mechanism test]\label{prop:qual}
The DAG $\bm X\to S\to L\to P$ with the moderating edge of $\bm X$ on the $S\to L$ link
makes \emph{process-tracing} predictions that are checkable case-by-case without numerical
estimation:
\begin{enumerate}
\item \emph{Mediation chain (mechanism)}: in a team where a structural intervention
(e.g.\ leader coaching) raises learning, one should observe an intervening rise in members'
expressed sense of safety \emph{temporally and narratively prior} to the increase in
feedback-seeking, error-discussion, etc. This is the qualitative signature of
$\bm X\to S\to L$ as opposed to a direct $\bm X\to L$ path with no safety change.
\item \emph{Moderation contrast}: comparing two teams with similar measured safety but
different structural support, the framework predicts that the \emph{same} level of safety
translates into more learning behavior in the better-supported team, because the slope is
$\gamma+\bm\delta^\top\bm x$. This is a difference observable in comparative case study.
\item \emph{Conditional independence}: the DAG implies $P\perp \bm X \mid L$ to the extent
$\bm\phi=\bm0$, and $L\perp\bm X\mid (S,\bm X S)$ residually; violations seen in case
narratives flag omitted paths and refine the model.
\end{enumerate}
These are exactly the implications a qualitative study (interviews, observation, archival
process data) can corroborate or refute, while Proposition~\ref{prop:quant} gives the
parallel quantitative tests.
\end{proposition}
\begin{proof}
Each item is a logical consequence of Model~\ref{mod:main}. Item~1: under the DAG, the only
paths from $\bm X$ to $L$ are the direct edge and the chain through $S$; if in a given team
the direct coefficient contribution $\bm\beta^\top\Delta\bm x$ is small relative to the
mediated contribution $(\gamma+\bm\delta^\top\bm x)\bm\alpha^\top\Delta\bm x$
(cf.\ \eqref{eq:decomp}), the temporal ordering $\bm X$-change $\to$ $S$-change $\to$
$L$-change is forced by the recursive structure $\bm X\prec S\prec L$. Item~2 is the
statement $\partial\mathbb E[L\mid\bm X,S]/\partial S=\gamma+\bm\delta^\top\bm X$ from
\eqref{eq:L}, evaluated at two $\bm x$ values. Item~3: in a recursive Gaussian (or, more
generally, linear) SEM the global Markov property holds on the DAG, so $d$-separation of $P$
and $\bm X$ given $L$ (which holds when the direct edge $\bm X\to P$ is absent,
$\bm\phi=\bm0$) yields the conditional independence $P\perp\bm X\mid L$; the residual
statement for $L$ follows because, conditional on its parents $(\bm X,S)$ and their product,
$L$ depends on $\bm X$ only through those regressors.
\end{proof}

\section{Conclusion}

Model~\ref{mod:main} is the requested integrated framework. Psychological safety enters
\emph{both} as a mediator (carrying the antecedent effect of structure onto learning,
Theorem~\ref{thm:decomp}) \emph{and} as a moderator (the interaction $(\bm\delta^\top\bm X)S$
conditions the safety--learning slope), and learning behavior is the proximal driver of
performance (Corollary~\ref{cor:perf}). Theorem~\ref{thm:ident} shows the framework is
identified and hence estimable, Theorem~\ref{thm:nonred} and Corollary~\ref{cor:bias} prove
that it strictly dominates any treatment of structure and safety as independent additive
predictors whenever moderation is present, and Propositions~\ref{prop:quant}--\ref{prop:qual}
give matched quantitative and qualitative test batteries. This jointly satisfies all
requirements of the problem statement.

\end{document}
